\chapter{Desarrollo Específico de la Contribución}

En este capítulo se describe, de forma concreta, el desarrollo técnico y analítico que constituye la aportación principal de este Trabajo de Fin de Máster en relación con el marco metodológico presentado anteriormente. El objetivo es fijar qué se ha implementado, en qué orden, y con qué criterios, de modo que la contribución quede claramente delimitada frente al estado del arte y la metodología general.

\section{Alcance y criterios de descripción}
Se explicitan el alcance del desarrollo (componentes, experimentos, métricas de interés) y las decisiones técnicas que definen la contribución frente a una mera replicación de arquitecturas de referencia. Inclúyase aquí la vinculación con los módulos VAE, GAN y criterio basado en la divergencia de Kullback-Leibler, según se haya materializado en el trabajo.

\section{Entorno experimental y despliegue}
Resúmase el entorno de software y hardware, versiones de librerías relevantes, y el procedimiento de entrenamiento y validación. Si existen fases (preprocesado, ajuste de hiperparámetros, prueba y evaluación de la cartera), consérvense en el orden lógico seguido en el desarrollo.

\section{Resultados intermedios y de validación}
Presentense los resultados de validación (por ejemplo, calidad del espacio latente, fidelidad de la \textit{Generative Adversarial Network}, estabilidad numérica) que sustenten la aportación antes de la sección de código y datos, si dicha entrega aplica a su estructura de memoria.
